\chapter{Einleitung}
\label{ch:chapter01}
Serverlose Anwendungen sind in den letzten Jahren immer beliebter geworden \cite{MarketsandMarkets2025ServerlessComputing}, da sie es Entwicklern ermöglichen,
sich auf die tatsächliche Geschäftslogik und die Implementierung der Anwendungen zu konzentrieren, während der 
Cloud-Anbieter die zugrunde liegende Infrastruktur zur Verfügung stellt und verwaltet. Mit der zunehmenden Verbreitung
und der wachsenden Komplexität von serverlosen Architekturen steigen auch die Anforderungen an die Observability dieser
Systeme. Distributed-Tracing-Frameworks haben sich hierzu als zentrale Lösung für das Observability-Problem 
solcher Systeme etabliert \cite{majors2026observability}. OpenTelemetry ist hierbei der De-facto-Standard für die Erfassung und den Export von 
Telemetriedaten \cite{AWSLambdaBestPractices}.

%
% Section: Motivation
%
\section{Motivation und Problemstellung}
\label{sec:intro:motivation}
Die Motivation für diese Arbeit ergibt sich aus drei Beobachtungen und basiert auf die Anregung der Autoren von 
\emph{\glqq A Comparison of Distributed Tracing Tools in Serverless Applications\grqq}\cite{10254754}. 
Erstens ist OpenTelemetry der De-facto-Standard für verteiltes Tracing in Cloud-nativen Anwendungen \cite{AWSLambdaBestPractices} und wird von 
vielen führenden Cloud-Anbietern nativ oder über dedizierte Distributionen unterstützt. Dazu zählen unter anderem 
\ac{AWS} \cite{AWSDistributedTracing}, Microsoft Azure \cite{AzureOpenTelemetry} und Google Cloud \cite{GoogleCloudOpenTelemetry}. Zweitens ist die Performance von 
verschiedenen Programmiersprachen im Kontext von \ac{FaaS} wie \ac{AWS} Lambda 
sehr variabel \cite{Rodriguez2024Evaluation}. Drittens verursacht die Instrumentierung von FaaS-Anwendungen mit OpenTelemetry einen nicht 
vernachlässigbaren Overhead in Bezug auf Laufzeit (CPU-Zeit), Speicherverbrauch und Kaltstartzeiten \cite{10254754}.
\\
Bisher existiert jedoch keine fundierte empirische Untersuchung, die den sprachspezifischen Overhead der 
OpenTelemetry-Instrumentierung in \ac{FaaS}-Umgebungen systematisch analysiert und quantifiziert. Daraus ergibt sich auch 
die zentrale Problemstellung dieser Arbeit – wie hoch ist der Overhead, der durch die Instrumentierung mit OpenTelemetry
in \ac{FaaS}-Umgebungen für Java, Python und Node.js verursacht wird, und wie variiert dieser Overhead zwischen den Sprachen?


\section{Forschungsfragen}
\label{sec:intro:fragen}

Aus der beschriebenen Problemstellung lassen sich folgende Forschungsfragen ableiten:

\paragraph{Leitfrage}
Wie unterscheidet sich der Performance-Overhead der Instrumentierung mit OpenTelemetry in AWS-Lambda-Anwendungen zwischen den Programmiersprachen Java, Python und Node.js, und welche Implikationen ergeben sich daraus für die Praxis?

\paragraph{Teilfragen}
\begin{enumerate}
    \item Wie unterscheidet sich der prozentuale Performance-Overhead (Laufzeit, Speicherverbrauch, Kaltstartzeit) zwischen Java, Python und Node.js?
    \item Bei welchen der fünf häufigsten Anwendungsfälle von \ac{FaaS} treten die größten sprachspezifischen Unterschiede auf?
    \item Welche praxisrelevanten Empfehlungen zur Sprach- und Konfigurationswahl lassen sich ableiten?
\end{enumerate}

% Section: Ziele
%
\section{Ziele der Arbeit}
\label{sec:intro:ziele}

Die Arbeit verfolgt drei operative Ziele:

\begin{itemize}
    \item \textbf{Z1 – Messung des Overheads:} Quantitative Bestimmung des Laufzeit-, Speicherverbrauch- und Kaltstartzeit-Overheads durch OpenTelemetry-Instrumentierung für Java, Python und Node.js auf AWS Lambda anhand der fünf häufigsten Anwendungsfälle von FaaS.
    
    \item \textbf{Z2 – Identifikation sprachspezifischer Unterschiede:} Systematischer Vergleich der fünf Anwendungsfälle (Microbenchmarks), 
    um diejenigen Szenarien zu identifizieren, in denen die größten sprachspezifischen Unterschiede auftreten.
    
    \item \textbf{Z3 – Praxisempfehlungen:} Ableitung konkreter Empfehlungen zur Sprachauswahl mit Hinblick auf den Performance-Overhead der Instrumentierung mit OpenTelemetry.
\end{itemize}

\section{Methodischer Ansatz}
\label{sec:intro:methode}

Die Arbeit folgt einem empirisch-quantitativen Ansatz mit experimentellem Design und gliedert sich in vier aufeinander aufbauende Phasen:

\begin{enumerate}
    \item \textbf{Vorbereitung und Grundlagen:} 
    Tiefgreifende Literaturrecherche zu Serverless Computing, \ac{FaaS}, Observability, verteiltem Tracing, OpenTelemetry und AWS X-Ray 
    sowie Analyse bestehender wissenschaftlicher Arbeiten zum Performance-Overhead von Tracing-Frameworks in Cloud-Umgebungen und zur Performance-Varianz von Programmiersprachen in \ac{FaaS}-Umgebungen.

    \item \textbf{Implementierung:} 
    Entwicklung von jeweils fünf Funktionen in Java, Python und Node.js, die die fünf häufigsten Anwendungsfälle von \ac{FaaS} repräsentieren.
    Für jede Funktion wird jeweils eine Version mit und ohne Instrumentierung mit OpenTelemetry erstellt, um den Overhead isoliert und direkt messen zu können.

    
    \item \textbf{Durchführung der Messungen:} 
    Ausführung aller Funktionen auf AWS Lambda unter kontrollierten Bedingungen (identische Speicherkonfiguration, Region, Wiederholungen usw.). 
    Erfassung der Performance-Metriken Laufzeit (CPU-Zeit), Speicherverbrauch und Kaltstartzeiten über AWS CloudWatch.
    
    \item \textbf{Auswertung und Empfehlungen:} 
    Berechnung des relativen prozentualen Overheads pro Sprache und Anwendungsfall. Sprach- und anwendungsfallübergreifender Vergleich, 
    statistische Auswertung der Messungen und Ableitung konkreter praxisrelevanter Empfehlungen zur Sprachauswahl.
\end{enumerate}

%
% Section: Struktur der Arbeit
%
\section{Gliederung}
\label{sec:intro:structure}

Die Arbeit gliedert sich in acht Kapitel.
\paragraph{Kapitel \ref{ch:chapter01}} führt in die Thematik ein, erläutert die Motivation, formuliert die Forschungsfragen, definiert die Ziele und beschreibt den methodischen Ansatz sowie die Struktur der Arbeit.

\paragraph{Kapitel \ref{ch:chapter02}} führt in die notwendigen Grundlagen ein: Cloud Computing und Microservices-Architektur, Serverless Computing und FaaS, AWS Lambda, Observability, Distributed Tracing sowie OpenTelemetry mit seinen Sampling und der AWS-Integration.

\paragraph{Kapitel \ref{ch:chapter03}} gibt einen Überblick über den Stand der Forschung zum Overhead von verteiltem Tracing und zur Performance-Varianz von Programmiersprachen in Serverless-Umgebungen. Dabei wird auch die Forschungslücke präzisiert, die später in der Arbeit behandelt wird.

\paragraph{Kapitel \ref{ch:chapter04}} beschreibt detailliert die Methodik und den Versuchsaufbau. Es umfasst das Experimentdesign, die fünf gängigsten Anwendungsfälle von \ac{FaaS}, die Implementierung der Testfunktionen in den drei Sprachen, die Durchführung der Messungen sowie die statistische Auswertung der Ergebnisse.

\paragraph{Kapitel \ref{ch:chapter05}} präsentiert die Ergebnisse der Messungen, gegliedert nach Performance und Performance-Overhead.

\paragraph{Kapitel \ref{ch:chapter06}} diskutiert die Ergebnisse im Hinblick auf die Forschungsfragen, vergleicht sie gegebenenfalls mit verwandten Arbeiten 
und leitet praxisrelevante Empfehlungen zur Sprachwahl ab.

\paragraph{Kapitel \ref{ch:chapter07}} fasst die Arbeit zusammen, hebt den wissenschaftlichen Beitrag hervor, reflektiert die Limitationen und gibt einen Ausblick auf mögliche Fragestellungen für spätere Arbeiten.
